\documentclass[12pt]{article}
\usepackage[utf8]{inputenc}
\usepackage[T1]{fontenc}
\usepackage{hyperref}
\title{Kadu Language Resources for K-12 Curriculum Development}
\author{Compiled by Claude Sonnet 5.5 (claude-sonnet-5-5) for Andr\'as Kornai}
\date{}
\begin{document}
\maketitle

\section{Introduction}
Kadu (ISO 639-3: zkd; also Kantu, autonym Asak) is a Sino-Tibetan language of the Luish group, spoken in the Sagaing Region of Myanmar. Wikipedia lists three varieties: Settaw, Mawkhwin and Mawteik, the last nearly extinct. The Banmauk Township area is the main centre of use today. Most people now speak Burmese. Kadu is not the same as Kaduo, a Ngwi language of Yunnan, China. Only general reference pages turned up in the searches.

\section{Primary Educational Resources}
No primers, textbooks or school programmes for Kadu turned up in the searches.

\section{Linguistic and Reference Materials}
\begin{itemize}
\item Wikipedia, Kadu language: \url{https://en.wikipedia.org/wiki/Kadu_language}
\item Wikipedia, Kadu people: \url{https://en.wikipedia.org/wiki/Kadu_people}
\item Numeral system survey for Kadu (Max Planck Institute for Evolutionary Anthropology): \url{https://lingweb.eva.mpg.de/channumerals/Kadu.htm}
\item Wikitongues language page for Kadu: \url{https://wikitongues.org/languages/zkd/}
\end{itemize}

\section{Religious and Community Resources}
\begin{itemize}
\item Joshua Project, Kadu, Asak in Myanmar: \url{https://joshuaproject.net/people_groups/21307/BM}
\end{itemize}

\section{Digital and Audio Resources}
No audio, video or software for Kadu turned up in the searches.

\section{Curriculum Development Context}
Little published material exists for Kadu. A curriculum would need new documentation, an orthography and teaching materials, ideally developed with the community. Burmese is the language of schooling in Myanmar.
\end{document}